\chapter{The Classical Maximum Principle}

Write
\begin{equation}
\tag{3.1}
Lu=a^{ij}D_{ij}u+b^iD_i u+cu,
\end{equation}
where $[a^{ij}]$ is positive definite and $\lambda$ denotes its least
eigenvalue. If $\Lambda$ is its greatest eigenvalue, then
\begin{equation}
\tag{3.2}
0<\lambda(x)|\xi|^2\leq a^{ij}(x)\xi_i\xi_j
\leq\Lambda(x)|\xi|^2.
\end{equation}
The lower-order coefficients are assumed to satisfy
\begin{equation}
\tag{3.3}
\frac{|b^i(x)|}{\lambda(x)}\leq C,\qquad i=1,\ldots,n.
\end{equation}
A function with $Lu=0$, $Lu\geq0$, or $Lu\leq0$ is called a
solution, subsolution, or supersolution, respectively.

\section*{3.1. The Weak Maximum Principle}

\begin{theorem}[Weak maximum principle]\label{gt-ch3-wmp}
\source{gilbarg-trudinger:page-0044}
\dcref{ch3:3.1}\group{ch03-classical-maximum-principle}
Let $\Omega$ be bounded, let $c=0$, and suppose $L$ is elliptic with
$|b|/\lambda$ bounded. If
$u\in C^2(\Omega)\cap C^0(\overline\Omega)$, then
\begin{equation}
\tag{3.4}
Lu\geq0\ (\leq0)\quad\text{in }\Omega,\qquad c=0,
\end{equation}
implies
\begin{equation}
\tag{3.5}
\sup_\Omega u=\sup_{\partial\Omega}u
\quad\left(\inf_\Omega u=\inf_{\partial\Omega}u\right).
\end{equation}
If continuity up to the boundary is omitted, replace the boundary extremum
by
\begin{equation}
\tag{3.6}
\sup_\Omega u=\limsup_{x\to\partial\Omega}u(x)
\quad\left(\inf_\Omega u=\liminf_{x\to\partial\Omega}u(x)\right).
\end{equation}
\end{theorem}
\begin{proof}
\sketch A function $w$ with $Lw>0$ cannot have an interior maximum, since
$Dw=0$ and $D^2w\leq0$ there. For sufficiently large $\gamma$,
\[
L(e^{\gamma x_1})\geq
\lambda(\gamma^2-\gamma\sup|b|/\lambda)e^{\gamma x_1}>0.
\]
Apply the strict case to $u+\varepsilon e^{\gamma x_1}$ and let
$\varepsilon\downarrow0$. The minimum assertion follows by replacing $u$
with $-u$.
\end{proof}

For $u^+=\max(u,0)$ and $u^-=\min(u,0)$, the sign condition $c\leq0$
gives the following extension.

\begin{corollary}\label{gt-ch3-cor-weak}
\source{gilbarg-trudinger:page-0045}
\dcref{ch3:3.2}\group{ch03-classical-maximum-principle}
Let $L$ be elliptic in bounded $\Omega$, let $c\leq0$, and let
$u\in C^2(\Omega)\cap C^0(\overline\Omega)$. Then
\begin{equation}
\tag{3.7}
Lu\geq0\quad(\leq0),\qquad c\leq0,
\end{equation}
implies
\begin{equation}
\tag{3.8}
\sup_\Omega u\leq\sup_{\partial\Omega}u^+
\quad\left(\inf_\Omega u\geq\inf_{\partial\Omega}u^-\right).
\end{equation}
If $Lu=0$, then
\begin{equation}
\tag{3.9}
\sup_\Omega|u|=\sup_{\partial\Omega}|u|.
\end{equation}
\end{corollary}
\begin{proof}
\sketch On $\Omega^+=\{u>0\}$, the operator with zero-order term removed
satisfies $L_0u=Lu-cu\geq0$. Apply Theorem~\ref{gt-ch3-wmp} on each
positive component. The remaining assertions follow by changing signs.
\end{proof}

\begin{theorem}[Comparison and uniqueness]\label{gt-ch3-comparison}
\source{gilbarg-trudinger:page-0045}
\dcref{ch3:3.3}\group{ch03-classical-maximum-principle}
Let $c\leq0$, and let $u,v\in C^2(\Omega)\cap C^0(\overline\Omega)$. If
$Lu\geq Lv$ in $\Omega$ and $u\leq v$ on $\partial\Omega$, then $u\leq v$
in $\Omega$. In particular, equal equations and equal boundary values imply
$u=v$.
\end{theorem}
\begin{proof}
\sketch Apply Corollary~\ref{gt-ch3-cor-weak} to $u-v$.
\end{proof}

\section*{3.2. The Strong Maximum Principle}

The boundary has an \emph{interior sphere} at $x_0$ if some ball
$B\subset\Omega$ has $x_0\in\partial B$.

\begin{lemma}[Hopf boundary point lemma]\label{gt-ch3-hopf}
\source{gilbarg-trudinger:page-0046}
\dcref{ch3:3.4}\group{ch03-classical-maximum-principle}
Suppose $L$ is uniformly elliptic, $c=0$, and $Lu\geq0$ in $\Omega$.
Assume that $u$ is continuous at $x_0\in\partial\Omega$, that
$u(x_0)>u(x)$ for $x\in\Omega$, and that $\Omega$ has an interior sphere at
$x_0$. If the outer normal derivative exists, then
\begin{equation}
\tag{3.10}
\frac{\partial u}{\partial\nu}(x_0)>0.
\end{equation}
If $c\leq0$ and $c/\lambda$ is bounded, the same holds when $u(x_0)\geq0$;
when $u(x_0)=0$, no sign condition on $c$ is needed. More generally, within
each fixed non-tangential cone at $x_0$,
\begin{equation}
\tag{3.11}
\liminf_{x\to x_0}
\frac{u(x_0)-u(x)}{|x-x_0|}>0.
\end{equation}
\end{lemma}
\begin{proof}
\sketch If $B_R(y)\subset\Omega$ touches $x_0$, use on
$B_R(y)\setminus\overline{B_\rho(y)}$ the barrier
$v=e^{-\alpha|x-y|^2}-e^{-\alpha R^2}$. Uniform ellipticity and the
coefficient bounds make $Lv\geq0$ for large $\alpha$. The weak maximum
principle applied to $u-u(x_0)+\varepsilon v$ gives
$\partial_\nu u(x_0)\geq-\varepsilon\partial_\nu v(x_0)>0$ and also the
non-tangential estimate. If $u(x_0)=0$, replace $L$ by $L-c^+$.
\end{proof}

\begin{theorem}[Strong maximum principle]\label{gt-ch3-smp}
\source{gilbarg-trudinger:page-0047}
\dcref{ch3:3.5}\group{ch03-classical-maximum-principle}
Let $L$ be uniformly elliptic on a domain $\Omega$.
If $c=0$ and $Lu\geq0$ (respectively $Lu\leq0$), an interior maximum
(respectively minimum) forces $u$ to be constant. If $c\leq0$ and
$c/\lambda$ is bounded, the same conclusion holds for a nonnegative maximum
or nonpositive minimum. Local uniform ellipticity and local coefficient
bounds suffice.
\end{theorem}
\begin{proof}
\sketch If a nonconstant $u$ attained its maximum $M$, take a largest ball
inside a component of $\{u<M\}$ that touches an interior point where $u=M$.
Lemma~\ref{gt-ch3-hopf} gives a nonzero derivative at that interior maximum,
a contradiction.
\end{proof}

\begin{theorem}[Neumann uniqueness]\label{gt-ch3-neumann}
\source{gilbarg-trudinger:page-0047}
\dcref{ch3:3.6}\group{ch03-classical-maximum-principle}
Let $\Omega$ be bounded and satisfy the interior sphere condition at every
boundary point. Suppose $L$ is uniformly elliptic, $c\leq0$, $c/\lambda$ is
bounded, and
$u\in C^2(\Omega)\cap C^0(\overline\Omega)$ satisfies $Lu=0$. If the normal
derivative exists and vanishes everywhere on $\partial\Omega$, then $u$ is
constant. If $c<0$ somewhere, then $u\equiv0$.
\end{theorem}
\begin{proof}
\sketch Unless $u$ is constant, the strong maximum principle makes either
$u$ or $-u$ attain a strict nonnegative maximum at the boundary. The Hopf
lemma contradicts the zero normal derivative. Substitution into $Lu=0$
forces a constant solution to vanish when $c<0$ somewhere.
\end{proof}

\section*{3.3. A Priori Bounds}

\begin{theorem}[Pointwise bound]\label{gt-ch3-apriori}
\source{gilbarg-trudinger:page-0048}
\dcref{ch3:3.7}\group{ch03-classical-maximum-principle}
Let $L$ be elliptic in bounded $\Omega$, let $c\leq0$, and put
$\beta=\sup_\Omega|b|/\lambda$. If
$u\in C^2(\Omega)\cap C^0(\overline\Omega)$ and $Lu\geq f$, then
\begin{equation}
\tag{3.12}
\sup_\Omega u\leq\sup_{\partial\Omega}u^+
+C\sup_\Omega\frac{|f^-|}{\lambda}.
\end{equation}
If $Lu=f$, then
\[
\sup_\Omega|u|\leq\sup_{\partial\Omega}|u|
+C\sup_\Omega\frac{|f|}{\lambda}.
\]
Here $C$ depends only on $\operatorname{diam}\Omega$ and $\beta$; if
$\Omega$ lies in a slab of width $d$, one may take
$C=e^{(\beta+1)d}-1$.
\end{theorem}
\begin{proof}
\sketch Place $\Omega$ in $0<x_1<d$. For $\alpha\geq\beta+1$,
$L_0e^{\alpha x_1}\geq\lambda$. Compare $u$ with
\[
v=\sup_{\partial\Omega}u^+
+(e^{\alpha d}-e^{\alpha x_1})
 \sup_\Omega\frac{|f^-|}{\lambda}.
\]
The weak maximum principle applied to $v-u$ proves (3.12); apply it also to
$-u$ when $Lu=f$.
\end{proof}

\begin{corollary}\label{gt-ch3-grad-cor}
\source{gilbarg-trudinger:page-0049}
\dcref{ch3:3.8}\group{ch03-classical-maximum-principle}
Let $Lu=f$ under the preceding regularity assumptions, without a sign
condition on $c$. If $C$ is the constant in Theorem~\ref{gt-ch3-apriori} and
\begin{equation}
\tag{3.13}
C_1=1-C\sup_\Omega\frac{c^+}{\lambda}>0,
\end{equation}
then
\begin{equation}
\tag{3.14}
\sup_\Omega|u|\leq\frac1{C_1}
\left(\sup_{\partial\Omega}|u|+C\sup_\Omega\frac{|f|}{\lambda}\right).
\end{equation}
\end{corollary}
\begin{proof}
\sketch Rewrite $(L_0+c)u=f$ as
$(L_0+c^-)u=f-c^+u$, apply Theorem~\ref{gt-ch3-apriori}, and absorb the
$C\sup(c^+/\lambda)\sup|u|$ term using (3.13).
\end{proof}

\section*{3.4. Gradient Estimates for Poisson's Equation}

\begin{theorem}[Interior gradient and log-Lipschitz estimates]\label{gt-ch3-poisson}
\source{gilbarg-trudinger:page-0053}
\dcref{ch3:3.9}\group{ch03-classical-maximum-principle}
Let $u\in C^2(\Omega)$ satisfy $\Delta u=f$. For
$d_x=\dist(x,\partial\Omega)$ and $d_{x,y}=\min(d_x,d_y)$,
if $Q=\{x:|x_i|<d\}$, then
\begin{equation}
\tag{3.15}
|D_i u(0)|\leq\frac nd\sup_{\partial Q}|u|
+\frac d2\sup_Q|f|,\qquad i=1,\ldots,n.
\end{equation}
Consequently,
\begin{equation}
\tag{3.16}
\sup_\Omega d_x|Du(x)|\leq C
\left(\sup_\Omega|u|+\sup_\Omega d_x^2|f(x)|\right),
\end{equation}
and, for $x\neq y$,
\begin{equation}
\tag{3.20}
d_{x,y}^2\frac{|Du(x)-Du(y)|}{|x-y|}
\leq C\left(\sup_\Omega|u|+\sup_\Omega d_x^2|f(x)|\right)
\left(1+\log\frac{d_{x,y}}{|x-y|}\right),
\end{equation}
with $C=C(n)$, interpreting the last factor by the corresponding coarse
bound when $|x-y|$ is comparable to $d_{x,y}$.
\end{theorem}
\begin{proof}
\sketch Odd reflection across the midplane of a cube and a quadratic
comparison function give (3.15); applying it on cubes of scale $d_x$ gives
(3.16). For the modulus estimate, set $M=\sup_Q|u|$, $N=\sup_Q|f|$, and
let $\mu$ be the gradient bound from (3.16). A double reflection in variables
$y,z$ is controlled by
\begin{equation}
\tag{3.17}
\psi(x',y,z)=\frac{4M|x'|^2}{d^2}+\frac{4\mu}{d}yz
+kyz\log\frac{2d}{y+z},
\end{equation}
where $k\geq\frac43(N+8(n-1)M/d^2)$. The maximum principle gives
\begin{equation}
\tag{3.18}
\frac12|u_y(0,y)-u_y(0,-y)|
\leq\frac{4\mu}{d}y+ky\log\frac{2d}{y}.
\end{equation}
The analogous tangential reflection gives
\begin{equation}
\tag{3.19}
\frac12|D_{n-1}u(0,z)-D_{n-1}u(0,-z)|
\leq\frac{4\mu}{d}z+\bar k z\log\frac{2d}{z},
\end{equation}
and the same estimate holds for $D_i$, $i\leq n-2$. Rotate coordinates along
the segment from $x$ to $y$ and rescale these estimates to obtain (3.20).
\end{proof}

\section*{3.5. A Harnack Inequality in Two Dimensions}

\begin{theorem}\label{gt-ch3-harnack}
\source{gilbarg-trudinger:page-0053, gilbarg-trudinger:page-0054, gilbarg-trudinger:page-0055, gilbarg-trudinger:page-0056}
\source{gilbarg-trudinger:page-0054}
\dcref{ch3:3.10}\group{ch03-classical-maximum-principle}
Let $u\geq0$ be a $C^2$ solution of
$au_{xx}+2bu_{xy}+cu_{yy}=0$ in the disk $D_R$, and suppose the operator is
uniformly elliptic. There is $K>0$, depending only on
$\mu=\sup\Lambda/\lambda$, such that
\begin{equation}
\tag{3.21}
Ku(0)\leq u(z)\leq K^{-1}u(0)\qquad(z\in D_{R/4}).
\end{equation}
\end{theorem}
\begin{proof}
\sketch Scale to $R=1$. Exponentials of two transverse parabolic defining
functions $v_\pm$ are subsolutions when their exponent depends only on $\mu$.
Writing $E_\pm=e^{\alpha v_\pm}$, normalize them as
\begin{equation}
\tag{3.22}
w_\pm=\frac{E_\pm-1}{e^{7\alpha/4}-1}.
\end{equation}
Comparison propagates a fixed fraction of $u(0)$ from the component of
$\{u>u(0)/2\}$ containing $0$ to a horizontal segment:
\begin{equation}
\tag{3.23}
u(x,1/2)>K_1u(0)\qquad(|x|\leq1/2).
\end{equation}
A second parabolic barrier gives
\begin{equation}
\tag{3.24}
u(z)>K_1K_2u(0)=Ku(0)\qquad(z\in D_{1/3}).
\end{equation}
Repeating the estimate in disks centered
at points of $D_{1/4}$ gives both inequalities in (3.21).
\end{proof}

The disk estimate also gives
\begin{equation}
\tag{3.25}
\sup_{D_{R/4}}u\leq\kappa\inf_{D_{R/4}}u.
\end{equation}

\begin{corollary}\label{gt-ch3-harnack-cor}
\source{gilbarg-trudinger:page-0056}
\dcref{ch3:3.11}\group{ch03-classical-maximum-principle}
Under the hypotheses of Theorem~\ref{gt-ch3-harnack}, for every
$\Omega'\Subset\Omega\subset\mathbb R^2$ there is
$\kappa=\kappa(\Omega,\Omega',\mu)$ such that
\begin{equation}
\tag{3.26}
\sup_{\Omega'}u\leq\kappa\inf_{\Omega'}u.
\end{equation}
\end{corollary}
\begin{proof}
\sketch Join points of $\Omega'$ by a finite chain of disks compactly
contained in $\Omega$ and iterate Theorem~\ref{gt-ch3-harnack}.
\end{proof}

The disk argument also applies to
\begin{equation}
\tag{3.27}
Lu=a^{ij}D_{ij}u+b^iD_i u+cu=0,
\qquad c\leq0,quad i,j=1,2,
\end{equation}
with constants depending additionally on the coefficient bounds and radius.

\begin{corollary}[Liouville theorem]\label{gt-ch3-liouville}
\source{gilbarg-trudinger:page-0056, gilbarg-trudinger:page-0057}
\dcref{ch3:3.12}\group{ch03-classical-maximum-principle}
An entire solution of a uniformly elliptic equation
$au_{xx}+2bu_{xy}+cu_{yy}=0$ on $\mathbb R^2$ that is bounded on one side is
constant.
\end{corollary}
\begin{proof}
\sketch Subtract the infimum and choose $z_0$ with $u(z_0)<\varepsilon$.
The disk estimate (3.21), whose constant is scale independent, bounds $u$ by
$K\varepsilon$ on every fixed point after taking a sufficiently large disk.
Let $\varepsilon\downarrow0$.
\end{proof}

\section*{3.6. Operators in Divergence Form}

For bounded measurable $a^{ij}$, put
\begin{equation}
\tag{3.28}
Lu=D_j(a^{ij}D_i u).
\end{equation}
A $C^1$ function is a generalized solution, subsolution, or supersolution
according as
\begin{equation}
\tag{3.29}
\int_\Omega a^{ij}D_i uD_j\phi\,dx=0\quad(\leq0,\ \geq0)
\end{equation}
for every nonnegative $\phi\in C_0^1(\Omega)$. Testing the subsolution
inequality
\begin{equation}
\tag{3.30}
\int_\Omega a^{ij}D_i uD_j\phi\,dx\leq0
\qquad(0\leq\phi\in C_0^1(\Omega))
\end{equation}
with the positive part of $u$ above its boundary supremum gives
\[
\int a^{ij}D_i uD_j u\,dx\leq0,
\]
and ellipticity yields the weak maximum principle without differentiating
the coefficients.
